\clearpage
\fancyhead[L]{\small JUNO.}
\hypertarget{native-268-lower}{}
\begin{center}
JUNO.
\end{center}

\noindent\rule{\linewidth}{0.4pt}

\begin{center}
Monatliche Correspondenz. Band XI. S. 475–482. Mai 1805.
\end{center}

\noindent\rule{\linewidth}{0.4pt}

Wir geben hier einen kleinen Nachtrag mit Beobachtungen der \emph{Juno}, die an weiter entfernten Sternwarten gemacht wurden und von denen wir erst später erfuhren. Sie umfassen den letzten Abschnitt ihrer diesmaligen Sichtbarkeit. Einige konnten bereits nicht mehr im Meridian gemacht werden und waren nur außerhalb des Meridians mit Äquatorialinstrumenten möglich. Vor allem wegen des hellen Mondscheins in den ersten Februartagen war \emph{Juno} nicht mehr zu sehen und ließ sich anschließend nicht wiederfinden. Nur Dr. \textsc{Gauss} hatte am 20. Februar noch das Glück, Juno zu beobachten. Es ist die letzte Beobachtung der Juno, die uns überhaupt mitgeteilt wurde.

Von \textsc{Piazzi} in Palermo erhielten wir noch eine Reihe von Meridianbeobachtungen, die er mit seinem Vollkreis gemacht hatte. Sie folgen hier.
\SourceLeaf{269}
\begin{center}
\begingroup\fontsize{9}{11}\selectfont\setlength{\tabcolsep}{3pt}\renewcommand{\arraystretch}{1.1}
\begin{tabular}{r|r|r|r}
1805 & \shortstack{Mittlere Zeit\\in Palermo} & \shortstack{Scheinbare\\Rektaszension der $\JunoSymbol$} & \shortstack{Scheinbare südliche\\Deklination der $\JunoSymbol$}\\ \hline
22. Nov. & $7^{\mathrm u}\ 42^{\mathrm m}\ 39^{\mathrm s}9$ & $357^\circ\ 12'\ 28''5$ & $.\ .\ .\ .\ .$ \\
24 & $7\ 36\ 12.5$ & $357\ 33\ 36.6$ & $10^\circ\ 49'\ 34''9$ \\
25 & $7\ 33\ 1.5$ & $357\ 44\ 54.9$ & $10\ 47\ 21.5$ \\
26 & $7\ 29\ 52.3$ & $357\ 56\ 36.9$ & $10\ 44\ 49.9$ \\
28 & $7\ 23\ 37.0$ & $358\ 20\ 49.5$ & $10\ 38\ 59.5$ \\
29 & $7\ 20\ 31.4$ & $358\ 33\ 26.1$ & $10\ 35\ 43.2$ \\
4. Dez. & $7\ 5\ 28.2$ & $359\ 42\ 42.3$ & $10\ 15\ 52.4$ \\
10 & $6\ 47\ 21.8$ & $1\ 5\ 11.9$ & $.\ .\ .\ .\ .$ \\
\end{tabular}\endgroup
\end{center}

\par\indent Die Beobachtung vom 10. Dezember scheint einen Fehler zu enthalten, da sie erheblich von der aus den Elementen berechneten Position abweicht. Dr. \textsc{Gauss} verglich diese Beobachtungen sorgfältig mit seinen IV. Elementen und erhielt die folgenden Ergebnisse.

\begin{center}
\begingroup\fontsize{9}{11}\selectfont\setlength{\tabcolsep}{3pt}\renewcommand{\arraystretch}{1.1}
\begin{tabular}{r|r|r|r|r}
1805 & \shortstack{Berechnete\\Rektaszension\\der $\JunoSymbol$} & \shortstack{Berechnete südliche\\Deklination\\der $\JunoSymbol$} & \multicolumn{2}{c}{\shortstack{Fehler der IV. Elemente\\von Dr. \textsc{Gauss}\\in Rekt.\qquad in Dekl.}}\\ \hline
22. Nov. & $357^\circ\ 12'\ 11''6$ & $.\ .\ .\ .\ .$ & $-16''9$ & $.\ .\ .$ \\
24 & $357\ 33\ 28.3$ & $10^\circ\ 49'\ 47''3$ & $-8.3$ & $+12''4$ \\
25 & $357\ 44\ 42.4$ & $10\ 47\ 25.9$ & $-12.5$ & $+4.4$ \\
26 & $357\ 56\ 18.3$ & $10\ 44\ 55.0$ & $-18.6$ & $+5.1$ \\
28 & $358\ 20\ 39.6$ & $10\ 39\ 5.0$ & $-9.9$ & $+5.5$ \\
29 & $358\ 33\ 22.3$ & $10\ 35\ 55.8$ & $-3,8$ & $+12.6$ \\
4. Dez. & $359\ 42\ 27.6$ & $10\ 16\ 9.3$ & $-14.7$ & $+16.9$ \\
10 &  &  &  &  \\
\end{tabular}\endgroup
\end{center}

\Needspace{14\baselineskip}
\par\indent Diese Abweichungen sind so gering, dass sie kaum eine Verbesserung der IV. Elemente der Juno zu erfordern scheinen. Denn alle unsere älteren Planetentafeln zeigen in der Rektaszension ebenso große, vielleicht oft noch größere Abweichungen. Dennoch stützte der fleißige Dr. \textsc{Gauss} die folgenden verbesserten V. Elemente der Juno auf diese Beobachtungen und auf eine eigene vom 20. Februar.

\begin{center}
\begingroup\fontsize{8}{10}\selectfont\setlength{\tabcolsep}{3pt}\renewcommand{\arraystretch}{1.1}
\begin{tabular}{lr}
Epoche 1805, Seeberger Meridian . . . . & $42^\circ\ 32'\ 36''$ \\
\multicolumn{2}{l}{$\left.\begin{array}{lr}\text{\AscendingNode\ 1805 . . . . . . . . . . . . . .}&171\ 4\ 15.6\\\text{Aphel . . . . . . . . . . . . .}&233\ 11\ 39\end{array}\right\}\begin{array}{l}\text{als siderisch ruhend}\\\text{vorausgesetzt}\end{array}$} \\
Neigung der Bahn . . . . . . . . . . & $13\ 3\ 38$ \\
Tägliche tropische Bewegung . . . . . . . & $815''9595$ \\
Jährliche . . . . . . . . . . . . . . & $82^\circ\ 43'\ 45''2$ \\
Exzentrizität . . . . . . . . . . . . & $0.254236$ \\
Logarithmus der Halbachse . . . . . . & $0.4256078$ \\
\end{tabular}\endgroup
\end{center}

\par\indent Die erwähnte Beobachtung von Dr. \textsc{Gauss} lautet wie folgt.

\begin{center}
\begingroup\fontsize{9}{11}\selectfont\setlength{\tabcolsep}{3pt}\renewcommand{\arraystretch}{1.1}
\begin{tabular}{c|c|c}
Mittlere Zeit & Rekt. $\JunoSymbol$ & Decl. bor. $\JunoSymbol$\\ \hline
1805, 20. Feb. $7^{\mathrm u}\ 11^{\mathrm m}\ 12^{\mathrm s}$ & $30^\circ\ 27'\ 2''$ & $1^\circ\ 47'\ 48''$ \\
\end{tabular}\endgroup
\end{center}

\par\indent Die aus den IV. Elementen für diese Zeit berechnete Position der \emph{Juno} lautet wie folgt.

\begin{center}
\begingroup\fontsize{9}{11}\selectfont\setlength{\tabcolsep}{3pt}\renewcommand{\arraystretch}{1.1}
\begin{tabular}{c|c|c|c}
& & \multicolumn{2}{c}{Unterschied}\\ Rekt. $\JunoSymbol$ & Decl. bor. $\JunoSymbol$ & in Rekt. & in Dekl.\\ \hline
$30^\circ\ 27'\ 41''6$ & $1^\circ\ 47'\ 20''$ & $+39''6$ & $-27''7$ \\
\end{tabular}\endgroup
\end{center}
\SourceLeaf{270}
\par\indent Aus diesen fünften Elementen berechnete Dr. \textsc{Gauss} die folgende Ephemeride für den geozentrischen Lauf der Juno [vom 20. Oktober 1805 bis zum 23. Juli 1806, wie für Ceres und Pallas auf S. \hyperlink{native-267}{262} und \hyperlink{native-268}{263} dieses Bandes].

[Im übrigen Teil dieses Aufsatzes gibt von \textsc{Zach} noch Beobachtungen von \textsc{Oriani} in Mailand und von \textsc{Olbers} in Bremen an und vergleicht sie mit den IV. Elementen von \textsc{Gauss}. Außerdem gibt er die von Dr. \textsc{Burckhardt} berechneten Elemente der Junobahn nach dem \emph{Moniteur} vom 28. Dezember 1804 wieder.]

[\textsc{Bode} veröffentlichte die auf Berliner Zeit bezogenen, von \textsc{Gauss} berechneten Elemente, die X\textsuperscript{ten} der Ceres, die VIII\textsuperscript{ten} der Pallas und die V\textsuperscript{ten} der Juno, im \emph{Astronomischen Jahrbuch für 1808}, Berlin 1805, \href{https://www.digitale-sammlungen.de/en/view/bsb10538337?page=280}{S. 270}.]

